\section*{Chapter \ref{ch:b3:organometallic-bonding} --- Organometallic Chemistry: Bonding and Ligands}

\begin{solution}{exo:b3:organometallic-bonding:1}
$\ce{Fe(CO)5}$: $8 + 10 = 18$. $\ce{Mn2(CO)10}$: $7 + 10 + 1 = 18$ per Mn. $\ce{CpMo(CO)3H}$: $6 + 5 + 6 + 1 =
18$. $\ce{[PtCl3(C2H4)]-}$: $10 + 3 + 2 + 1 = 16$. $\ce{Cr($\eta^6$-C6H6)(CO)3}$: $6 + 6 + 6 = 18$. $\ce{Cp2ZrCl2}$:
$4 + 10 + 2 = 16$.
\end{solution}

\begin{solution}{exo:b3:organometallic-bonding:2}
Fe(0) $d^8$; Mn(0) $d^7$; Mo(II) $d^4$; Pt(II) $d^8$; Cr(0) $d^6$; Zr(IV) $d^0$.
\end{solution}

\begin{solution}{exo:b3:organometallic-bonding:3}
$\ce{[Mn(CO)6]+} > \ce{Cr(CO)6} > \ce{[V(CO)6]-}$: the more electron-rich the metal, the more it
back-donates into $\pi^*$ and the weaker the C--O bonds.
\end{solution}

\begin{solution}{exo:b3:organometallic-bonding:4}
Ethane ($\ce{Mn(CO)5}$ and $\ce{CpFe(CO)2}$ are \omterm{def:b3:organometallic-bonding:isolobal}{isolobal} with \ce{CH3}); ethane again; tetrahedrane
$\ce{(CH)4}$ with three CH replaced by $\ce{Co(CO)3}$.
\end{solution}

\begin{solution}{exo:b3:organometallic-bonding:5}
$fac$ ($C_{3v}$): two (A$_1$ + E). $mer$ ($C_{2v}$): three (2A$_1$ + B$_1$).
\end{solution}

\begin{solution}{exo:b3:organometallic-bonding:6}
The bulky phosphine crowds the metal: dissociation relieves the strain, so it is
faster, and low-coordinate intermediates are favoured.
\end{solution}

\begin{solution}{exo:b3:organometallic-bonding:7}
The chromium complex is a \omterm{def:b3:organometallic-bonding:carbene}{Fischer carbene} (OMe substituent, Cr(0), CO ligands): an
amine attacks the electrophilic \omterm{def:b3:organometallic-bonding:carbene}{carbene} carbon and replaces OMe. The tantalum
alkylidene is a \omterm{def:b3:organometallic-bonding:carbene}{Schrock carbene}: its nucleophilic carbon attacks the carbonyl of a
ketone, giving an alkene and a Ta=O unit, as in a Wittig reaction.
\end{solution}

\begin{solution}{exo:b3:organometallic-bonding:8}
Staggered: no $\delta$ overlap, the two $\delta$ electrons are non-bonding: bond order 3. $\ce{[Re2Cl8]^{4-}}$:
$\sigma^2\pi^4\delta^2\delta^{\ast2}$, bond order 3.
\end{solution}

\begin{solution}{exo:b3:organometallic-bonding:9}
A short M$\cdots$H distance (diffraction, best neutron); a low C--H stretching
wavenumber; an NMR signal at high field with a reduced one-bond C--H coupling.
\end{solution}

\begin{solution}{exo:b3:organometallic-bonding:10}
16: square-planar $d^8$ ($\ce{[PtCl4]^{2-}}$, Wilkinson's catalyst), whose 18th-electron orbital
is too high; early-metal \omterm{def:b3:organometallic-bonding:metallocene}{metallocenes} such as $\ce{Cp2ZrCl2}$, too crowded for more
ligands. 17: $\ce{V(CO)6}$, which cannot dimerise for steric reasons. 19: cobaltocene,
whose extra electron sits in an antibonding orbital (and is easily lost).
\end{solution}

\begin{solution}{exo:b3:organometallic-bonding:11}
CO is a strong $\pi$ acceptor: it lowers the $t_{2g}$ level and makes $\Delta_{\mathrm o}$ large, far above the
pairing energy: $t_{2g}^6$, diamagnetic. The $d$--$d$ transitions, at $\Delta_{\mathrm o}$, fall in the ultraviolet,
and no visible light is absorbed.
\end{solution}

\begin{solution}{exo:b3:organometallic-bonding:12}
Donation from $\pi$ and back-donation into $\pi^*$ both lower the C--C bond order. In Zeise's
salt (Pt(II), moderate back-donation) the bond lengthens a little; a low-valent,
electron-rich metal back-donates strongly and the complex approaches the
\omterm{def:b3:organometallic-bonding:dcd}{metallacyclopropane} limit, with a C--C bond close to a single bond.
\end{solution}

\begin{solution}{pb:b3:organometallic-bonding:1}
\textbf{1.} $\ce{Fe^{2+}}$ ($d^6$) $+ 2 \times 6 = 18$.
\textbf{2.} $8 + 2 \times 5 = 18$.
\textbf{3.} $9 + 10 = 19$.
\textbf{4.} $10 + 10 = 20$.
\textbf{5.} $18 - 1 = 17$.
\textbf{6.} Fe(II) $d^6$; Co(II) $d^7$; Ni(II) $d^8$; Fe(III) $d^5$.
\textbf{7.} $a_{1g}$ ($d_{z^2}$), $e_{1g}$ ($d_{xz}$, $d_{yz}$), $e_{2g}$ ($d_{xy}$, $d_{x^2-y^2}$).
\textbf{8.} The $e_{1g}$ pair overlaps strongly with the filled ring $e_{1g}$ combination: bonding
orbitals mostly on the rings, antibonding $e_{1g}^\ast$ mostly on the metal, high in
energy.
\textbf{9.} $e_{2g} \approx a_{1g} < e_{1g}^\ast$.
\textbf{10.} $e_{2g}^4a_{1g}^2$: no unpaired electron.
\textbf{11.} One electron in $e_{1g}^\ast$: one unpaired electron, $1.73\,\mu_B$.
\textbf{12.} Two electrons in the two $e_{1g}^\ast$ orbitals, parallel: two unpaired, $\sqrt{8} =
2.83\,\mu_B$.
\textbf{13.} $e_{2g}^3a_{1g}^2$: one unpaired electron. An $e_{2g}^3$ hole gives an orbitally degenerate
(E) state with an \omterm{def:b3:complex-spectra-magnetism:moment}{orbital contribution}.
\textbf{14.} Its 19th electron is in an antibonding orbital and is easily removed,
giving the 18-electron cobaltocenium cation.
\textbf{15.} The electron removed comes from a nearly non-bonding orbital: the
structure hardly changes (small \omterm{def:b3:complex-mechanisms:reorganisation}{reorganisation energy}, \cref{ch:b3:complex-mechanisms}).
\textbf{16.} The couple is fast and reversible in most solvents, the compound is
soluble and stable, and its potential depends little on the solvent because the
charge is spread over a large ion.
\textbf{17.} With 20 electrons, two of them antibonding, it reacts by losing them: it
adds ligands with slippage or loss of a ring, or is oxidised.
\textbf{18.} $\ce{CpFe(CO)2}$: $8 + 5 + 4 = 17$ electrons, one short of 18, one frontier orbital: \omterm{def:b3:organometallic-bonding:isolobal}{isolobal}
with \ce{CH3}. The dimer $\ce{[CpFe(CO)2]2}$ is ``ethane'', with an Fe--Fe bond (in practice two CO
ligands bridge it).
\textbf{19.} $\ce{CpFe(CO)2-Mn(CO)5}$, with an Fe--Mn bond.
\textbf{20.} $\ce{PPh3}$ is a better donor and a poorer $\pi$ acceptor than CO: the metal back-donates
more to the remaining CO, whose stretches fall.
\textbf{21.} Two (A$_1$ + E).
\textbf{22.} Ring slippage to $\eta^3$ frees two electrons' worth of coordination: an entering
ligand can bind associatively without the metal exceeding 18 electrons.
\textbf{23.} \textbf{$\mu_{\text{spin only}}(\text{nickelocene}) = \sqrt{2 \times 4} = 2.83\,\mu_B$.}
\end{solution}
